\section{Instalación de Electricidad}

En esta unidad solo se recogen las particularidades de Electricidad. El uso general de la herramienta ya se ha explicado en los apartados anteriores, por lo que aquí se detallan únicamente las reglas, avisos y casos propios de esta instalación.

\subsection{Descripción general de Electricidad}

\begin{itemize}
\item Esta instalación sirve para dibujar y editar recorridos de electricidad dentro de Allplan.
\item Está pensada para trabajar tanto con \texttt{rejibands} como con conductos o corrugados asociados a distintos usos eléctricos.
\item Lo que la diferencia de otras instalaciones es la variedad de recorridos eléctricos específicos, la coexistencia de capas de corrugados y de telecomunicaciones, y la presencia de variantes \texttt{TD} en varios scripts.
\end{itemize}

\subsection{Sistemas disponibles en Electricidad}

En la configuración actual de Electricidad aparecen estos sistemas:

\begin{itemize}
\item \texttt{Rejiband U 100mm};
\item \texttt{Rejiband U 200mm};
\item \texttt{Telecomunicaciones};
\item \texttt{Luz retorno / Paralelas};
\item \texttt{Alimentación horno};
\item \texttt{Alimentación luces / Cajetines};
\item \texttt{Cajetín a enchufe};
\item \texttt{Interruptores / Domótica}.
\end{itemize}

Las diferencias principales entre ellos son estas:

\begin{itemize}
\item \texttt{Rejiband U 100mm} y \texttt{Rejiband U 200mm} trabajan como soporte lineal principal con ancho diferente.
\item El resto trabajan como conductos o corrugados asociados a usos concretos de la instalación eléctrica.
\item Cada tipo define color, diámetro base, solape y lógica geométrica propia.
\end{itemize}

\manualimage{CAPTURA-ELECTRICIDAD-01.png}{Paleta de Electricidad con los sistemas disponibles visibles}

\subsection{Distribución y variantes en Electricidad}

En la paleta existe el campo \texttt{Tipo de Distribucion} y en los scripts aparecen variantes específicas con sufijo \texttt{TD}, por ejemplo:

\begin{itemize}
\item \texttt{conducto\_telecomunicaciones\_td};
\item \texttt{conducto\_luz\_retorno\_paralelas\_td};
\item \texttt{conducto\_alimentacion\_horno\_td};
\item \texttt{conducto\_alimentacion\_luces\_cajetines\_td};
\item \texttt{conducto\_cajetin\_a\_enchufe\_td};
\item \texttt{conducto\_interruptores\_domotica\_td}.
\end{itemize}

Esto indica que Electricidad distingue al menos entre una lógica base y otra variante \texttt{TD}.

Además, la paleta incluye también el campo \texttt{Cara (EN)}, lo que sugiere otra lógica adicional ligada a la cara de colocación.

Por tanto:

\begin{itemize}
\item \pendingtag\ debe validarse cómo selecciona realmente el usuario la variante \texttt{TD} o \texttt{EN} en Electricidad;
\item \pendingtag\ debe comprobarse qué combinación práctica existe entre \texttt{Tipo de Distribucion} y \texttt{Cara (EN)}.
\end{itemize}

\subsection{Diámetros y cambios de diámetro en Electricidad}

\subsubsection{Diámetros por sistema}

Los diámetros confirmados en la configuración actual son:

\begin{itemize}
\item \texttt{Rejiband U 100mm}: \texttt{100 mm};
\item \texttt{Rejiband U 200mm}: \texttt{200 mm};
\item \texttt{Telecomunicaciones}: \texttt{20 mm};
\item \texttt{Luz retorno / Paralelas}: \texttt{20 mm};
\item \texttt{Alimentación horno}: \texttt{25 mm};
\item \texttt{Alimentación luces / Cajetines}: \texttt{20 mm};
\item \texttt{Cajetín a enchufe}: \texttt{20 mm};
\item \texttt{Interruptores / Domótica}: \texttt{20 mm}.
\end{itemize}

Esto significa que cada recorrido eléctrico queda ligado al diámetro definido por su sistema.

\manualimage{CAPTURA-ELECTRICIDAD-02.png}{Selector de diámetro en Electricidad mostrando diferentes sistemas y diámetros}

\subsubsection{Qué ocurre al modificar el diámetro}

La paleta incluye el botón \texttt{Modificar diámetro}.

En esta instalación, el cambio de diámetro debe revisarse con especial cuidado porque:

\begin{itemize}
\item cada sistema ya parte de un diámetro base concreto;
\item en \texttt{Rejiband} el ancho es parte esencial del resultado;
\item y en los conductos eléctricos el diámetro afecta a la geometría, al área y a varios atributos \texttt{pmp\_*}.
\end{itemize}

\subsubsection{Solape y continuidad}

En Electricidad, varios sistemas usan solape adicional (\texttt{overlap\_mm}) para mantener continuidad geométrica:

\begin{itemize}
\item \texttt{Rejiband U}: \texttt{50 mm};
\item la mayoría de conductos eléctricos: \texttt{10 mm};
\item \texttt{Alimentación horno}: \texttt{12.5 mm}.
\end{itemize}

Esto es importante porque el resultado final no depende solo de la longitud dibujada, sino también del solape definido por el sistema.

\subsection{Reglas geométricas y validaciones en Electricidad}

\subsubsection{Ángulos permitidos}

La configuración general de Electricidad indica:

\begin{itemize}
\item \texttt{Libre (ángulo interior >= 45°)}.
\end{itemize}

Esto la diferencia de instalaciones como Ventilación, donde el catálogo de ángulos está más restringido.

Por tanto, en Electricidad la regla práctica es:

\begin{itemize}
\item el recorrido puede ser libre;
\item pero el ángulo interior no debe bajar de \texttt{45 grados}.
\end{itemize}

\subsubsection{Longitud mínima}

La longitud mínima registrada para todos los sistemas principales de Electricidad es:

\begin{itemize}
\item \texttt{300 mm}.
\end{itemize}

\pendingtag\ El mensaje exacto mostrado cuando no se cumple la longitud mínima debe validarse en Allplan.

\subsubsection{Orientación 3D y tramos inclinados}

Electricidad utiliza la misma lógica general de orientación 3D compartida por la base del sistema, pero aquí influye especialmente en:

\begin{itemize}
\item rejibands;
\item conductos eléctricos;
\item y elementos definidos conectados al trazado.
\end{itemize}

Por eso, además de la planta, conviene revisar cómo queda la orientación final de los recorridos en cambios de plano o en tramos inclinados.

\manualimage{CAPTURA-ELECTRICIDAD-03.png}{Ejemplo de orientación 3D o tramo inclinado en Electricidad}

\subsection{Layers y atributos específicos de Electricidad}

\subsubsection{Layers base}

Los layers base confirmados en el registro de Electricidad son:

\begin{itemize}
\item \texttt{IS\_REJIBANDS};
\item \texttt{IS\_COR\_TELECOS\_FAB};
\item \texttt{IS\_COR\_TELECOS\_IN};
\item \texttt{IS\_COR\_TELECOS\_OUT};
\item \texttt{IS\_COR\_ELECTRICITAT\_FAB};
\item \texttt{IS\_COR\_ELECTRICITAT\_FAB\_IN};
\item \texttt{IS\_COR\_ELECTRICITAT\_FAB\_OUT};
\item \texttt{KN\_ELECTRICITAT};
\item \texttt{KN\_XPS\_RECESS};
\item \texttt{KN\_X\_ELECTRICITAT};
\item \texttt{KN\_Y\_ELECTRICITAT}.
\end{itemize}

Esto indica una separación clara entre:

\begin{itemize}
\item rejibands;
\item corrugados eléctricos;
\item telecomunicaciones;
\item y variantes relacionadas con \texttt{TD} o con cara \texttt{X} y \texttt{Y}.
\end{itemize}

\manualimage{CAPTURA-ELECTRICIDAD-04.png}{Ejemplo de layers aplicados sobre recorridos de Electricidad}

\subsubsection{Atributos detectados en scripts}

En los scripts de Electricidad se detectan, entre otros, estos atributos:

\begin{itemize}
\item \texttt{6\_CC\_IS};
\item \texttt{pmp\_altura};
\item \texttt{pmp\_amplada};
\item \texttt{pmp\_area};
\item \texttt{pmp\_CARTICULO};
\item \texttt{pmp\_color};
\item \texttt{pmp\_diametre};
\item \texttt{pmp\_longitud};
\item \texttt{pmp\_nom};
\item \texttt{pmp\_pare};
\item \texttt{pmp\_seccio};
\item y otros atributos \texttt{pmp\_*} de clasificación y propiedades físicas.
\end{itemize}

\subsubsection{Qué conviene revisar en los atributos}

En Electricidad, varios scripts rellenan atributos a partir del diámetro y del color del sistema.

Esto conviene revisarlo especialmente en:

\begin{itemize}
\item \texttt{Alimentación horno};
\item \texttt{Telecomunicaciones};
\item \texttt{Rejiband U};
\item y en las variantes \texttt{TD}.
\end{itemize}

\subsection{Accesorios y comportamiento geométrico en Electricidad}

Electricidad no registra un catálogo explícito de accesorios automáticos como sí ocurre en Agua o Ventilación.

Lo que sí se observa es esto:

\begin{itemize}
\item los recorridos eléctricos se generan como tramos continuos;
\item en varios casos se prioriza la continuidad geométrica por solape;
\item \texttt{Rejiband U} se comporta como geometría lineal propia;
\item la unión booleana o continuidad del recorrido es especialmente importante en esta instalación.
\end{itemize}

De hecho, en la configuración se deja indicado explícitamente que:

\begin{itemize}
\item trabajar en modo grupo es crítico para no romper la continuidad de ciertos recorridos de electricidad.
\end{itemize}

Esto afecta sobre todo a:

\begin{itemize}
\item \texttt{Telecomunicaciones};
\item \texttt{Luz retorno / Paralelas};
\item \texttt{Alimentación horno};
\item \texttt{Alimentación luces / Cajetines};
\item \texttt{Cajetín a enchufe};
\item \texttt{Interruptores / Domótica}.
\end{itemize}

\subsection{Elementos definidos en Electricidad}

En la información actual, el elemento definido visible y coherente con la paleta es:

\begin{itemize}
\item \texttt{Caixa Connexions 200}.
\end{itemize}

Este elemento aparece tanto en el registro de Electricidad como en la página \texttt{Elemento} del archivo \texttt{electricidad\_polyline.pyp}.

Por eso, a diferencia de Ventilación, aquí no hay incoherencia entre el registro y la paleta en este punto.

Su función práctica es insertar una caja de conexiones sobre un punto del recorrido.

\manualimage{CAPTURA-ELECTRICIDAD-05.png}{Uso de Caixa Connexions 200 en Electricidad}

\subsubsection{Tipos de punto admitidos}

En la paleta de Electricidad, \texttt{Caixa Connexions 200} admite:

\begin{itemize}
\item \texttt{Punto inicial};
\item \texttt{Punto final};
\item \texttt{Intermedio o libre}.
\end{itemize}

\subsection{Macros en Electricidad}

Electricidad incluye la misma lógica general de macros compartida por la infraestructura base:

\begin{itemize}
\item selección de punto;
\item selección de macro SmartSymbol;
\item cota manual o relativa a local;
\item confirmación de inserción.
\end{itemize}

Los campos visibles en paleta son coherentes con ese flujo.

\pendingtag\ si Electricidad usa macros con alguna regla específica adicional.

\manualimage{CAPTURA-ELECTRICIDAD-06.png}{Página de macros de Electricidad}

\subsection{Soportes en Electricidad}

Electricidad dispone también de página específica de soportes.

Los campos confirmados en paleta son:

\begin{itemize}
\item \texttt{Tipo de soporte};
\item \texttt{Subtipo instalación};
\item \texttt{Superficie};
\item \texttt{Cota A};
\item \texttt{Cota B};
\item \texttt{Ángulo de inclinación};
\item modos de edición de soportes;
\item y \texttt{Atributo de soporte}.
\end{itemize}

La estructura de esta página es equivalente a la de otras instalaciones.

Por eso, en Electricidad conviene revisar especialmente:

\begin{itemize}
\item si el subtipo visible se corresponde realmente con \texttt{Electricidad};
\item qué tipos de soporte están habilitados en entorno real;
\item y cómo se clasifica el soporte una vez creado.
\end{itemize}

\manualimage{CAPTURA-ELECTRICIDAD-07.png}{Bloque de soportes con subtipo Electricidad}

\subsection{Copias, duplicados y comportamiento especial en Electricidad}

Electricidad usa la misma lógica general de serialización, restauración, layers, atributos y posible copia relacionada que el resto de instalaciones basadas en \texttt{PolyLib}.

Además, hay un comportamiento técnico especialmente importante en esta instalación:

\begin{itemize}
\item en varios recorridos eléctricos la continuidad geométrica es crítica;
\item por eso el modo agrupado tiene más peso que en otros sistemas.
\end{itemize}

Esto significa que el usuario debe revisar no solo que el recorrido exista, sino que también se esté resolviendo como un conjunto continuo y no como tramos aislados cuando eso no corresponde.

\pendingtag\ si Electricidad está usando en producción copias automáticas a otros archivos de dibujo y en qué casos.

\subsection{Inconsistencias o puntos a validar en Electricidad}

Aquí hay dos puntos que conviene dejar documentados:

\begin{enumerate}
\item En el registro aparecen variantes \texttt{TD} que no se muestran directamente como tipo de instalación independiente en la lista principal.
\item Algunas etiquetas de \texttt{elements3D} parecen heredadas o poco coherentes, por ejemplo:
\begin{itemize}
\item \texttt{conducto\_cajetin\_a\_enchufe} aparece con la etiqueta \texttt{Conducto Aislado 150 mm};
\item \texttt{conducto\_interruptores\_domotica} aparece con la etiqueta \texttt{Conducto Aislado 160 mm}.
\end{itemize}
\end{enumerate}

Esto sugiere que parte de la definición visible puede venir de una reutilización de configuración y conviene validarla antes de cerrar el manual.

\subsection{Puntos a revisar con especial atención en Electricidad}

Antes de finalizar una instalación de Electricidad conviene revisar especialmente:

\begin{itemize}
\item que el sistema activo sea realmente el que corresponde al uso eléctrico que se quiere modelar;
\item que el diámetro base sea el correcto para ese sistema;
\item que los recorridos sobre \texttt{Rejiband} mantengan la continuidad esperada;
\item que las variantes \texttt{TD} o \texttt{EN} se estén aplicando como realmente necesita el proyecto;
\item que los layers elegidos respondan al criterio de telecomunicaciones, corrugados o cara de colocación;
\item que \texttt{Caixa Connexions 200} esté colocada en el punto adecuado;
\item que los soportes se hayan revisado con subtipo, superficie, cotas y tipo correctos;
\item y que las etiquetas visibles en paleta coincidan con el comportamiento real del sistema.
\end{itemize}

